\documentclass[11pt,letterpaper]{article}
\usepackage[T1]{fontenc}
\usepackage[utf8]{inputenc}
\usepackage[margin=0.88in,top=0.84in,bottom=0.83in,headheight=13pt,headsep=20pt,footskip=28pt]{geometry}
\usepackage{newpxtext}
\usepackage{amsmath,amsthm}
\usepackage{newpxmath}
% newpxtext selects TeX Gyre Heros (qhv), scaled to .94, for sans serif.
\usepackage{microtype}
\usepackage{xcolor}
\definecolor{Forest}{HTML}{30392C}
\definecolor{Olive}{HTML}{687144}
\definecolor{Muted}{HTML}{65685F}
\definecolor{Sage}{HTML}{CBD0BC}
\definecolor{Pale}{HTML}{F2F4EB}
\usepackage{mathtools}
\usepackage{booktabs,tabularx,longtable,array}
\usepackage{enumitem}
\usepackage{titlesec}
\usepackage{fancyhdr}
\usepackage[most]{tcolorbox}
\usepackage{needspace}
\usepackage{xurl}
\usepackage[colorlinks=true,linkcolor=Olive,citecolor=Olive,urlcolor=Olive,bookmarksnumbered=true,pdfencoding=auto,psdextra]{hyperref}
\hypersetup{pdftitle={Large cardinals at the inconsistency frontier: research continuation},pdfauthor={Prepared for Vladimir Reshetnikov},pdfsubject={A Prikry consistency refinement and local definability obstructions},pdfkeywords={large cardinals, cover exacting, Prikry forcing, rank into rank, HOD, relative consistency}}
\urlstyle{same}
\setlength{\parindent}{0pt}
\setlength{\parskip}{5.1pt plus 1pt minus .4pt}
\linespread{1.035}
\setlength{\emergencystretch}{2em}
\setcounter{tocdepth}{1}
\setcounter{secnumdepth}{2}
\titleformat{\section}{\Large\bfseries\color{Forest}}{\thesection}{.6em}{}
\titleformat{\subsection}{\large\bfseries\color{Forest}}{\thesubsection}{.55em}{}
\titleformat{\subsubsection}{\normalsize\bfseries\color{Forest}}{}{0pt}{}
\titlespacing*{\section}{0pt}{18pt plus 3pt minus 2pt}{8pt}
\titlespacing*{\subsection}{0pt}{12pt plus 2pt minus 1pt}{5pt}
\titlespacing*{\subsubsection}{0pt}{9pt}{3pt}
\setlist[itemize]{leftmargin=1.4em,itemsep=2pt,topsep=3pt,parsep=0pt}
\setlist[enumerate]{leftmargin=1.7em,itemsep=3pt,topsep=4pt,parsep=0pt}
\pagestyle{fancy}
\fancyhf{}
\fancyhead[L]{\sffamily\fontsize{9}{11}\selectfont\color{Muted}LARGE CARDINALS AT THE INCONSISTENCY FRONTIER}
\fancyhead[R]{\sffamily\fontsize{9}{11}\selectfont\color{Muted}18 SEP 2026}
\fancyfoot[L]{\small\color{Muted}Prikry consistency and definability obstructions}
\fancyfoot[R]{\small\color{Muted}\thepage}
\renewcommand{\headrulewidth}{.35pt}
\renewcommand{\footrulewidth}{0pt}
\fancypagestyle{plain}{\fancyhf{}\fancyfoot[L]{\small\color{Muted}Prikry consistency and definability obstructions}\fancyfoot[R]{\small\color{Muted}\thepage}\renewcommand{\headrulewidth}{0pt}}
\tcbset{colback=Pale,colframe=Sage,colbacktitle=Sage,coltitle=Forest,fonttitle=\bfseries,boxrule=.5pt,arc=3pt,left=9pt,right=9pt,top=6pt,bottom=6pt,toptitle=3pt,bottomtitle=3pt,before skip=8pt,after skip=8pt,breakable}
\newtcolorbox{assessment}[1]{title={#1}}
\theoremstyle{definition}
\newtheorem{definition}{Definition}[section]
\newtheorem{theorem}[definition]{Theorem}
\newtheorem{proposition}[definition]{Proposition}
\newtheorem{lemma}[definition]{Lemma}
\newtheorem{corollary}[definition]{Corollary}
\newcommand{\ZFC}{\mathsf{ZFC}}
\newcommand{\ZF}{\mathsf{ZF}}
\newcommand{\AC}{\mathsf{AC}}
\newcommand{\DC}{\mathsf{DC}}
\newcommand{\HOD}{\mathrm{HOD}}
\newcommand{\OD}{\mathrm{OD}}
\newcommand{\HH}{\mathsf{HH}}
\newcommand{\Ex}{\operatorname{Ex}}
\newcommand{\UEx}{\operatorname{UEx}}
\newcommand{\Ext}{\operatorname{Ext}}
\newcommand{\SC}{\operatorname{SC}}
\newcommand{\CEx}{\operatorname{CEx}}
\newcommand{\Con}{\operatorname{Con}}
\newcommand{\crit}{\operatorname{crit}}
\newcommand{\cf}{\operatorname{cf}}
\newcommand{\ot}{\operatorname{ot}}
\newcommand{\Ord}{\mathrm{Ord}}
\newcommand{\Card}{\mathrm{Card}}
\newcommand{\rank}{\operatorname{rank}}
\newcommand{\id}{\mathrm{id}}
\newcommand{\restr}{\mathbin{\upharpoonright}}
\newcommand{\Izero}{\mathrm{I0}}
\newcommand{\Ione}{\mathrm{I1}}
\newcommand{\Itwo}{\mathrm{I2}}
\newcommand{\Ithree}{\mathrm{I3}}
\newcommand{\Isharp}{\mathrm{I0}^{\#}}
\newcommand{\Status}[1]{\textbf{Status.} #1}
\newcommand{\Priority}[1]{\textbf{Research priority.} #1}
\newcolumntype{Y}{>{\raggedright\arraybackslash}X}
\newcolumntype{P}[1]{>{\raggedright\arraybackslash}p{#1}}
\newcommand{\arxiv}[1]{\href{https://arxiv.org/abs/#1}{arXiv:#1}}
\newcommand{\doi}[1]{\href{https://doi.org/#1}{doi:\nolinkurl{#1}}}

\newcommand{\Iwf}{\mathrm{I3}_{\mathrm{wf}(0)}}
\newcommand{\Iwfone}{\mathrm{I3}_{1}}
\newcommand{\RelEx}{\operatorname{RelEx}}
\newcommand{\Bad}{\operatorname{Bad}}
\newcommand{\VP}{\mathsf{VP}}
\newcommand{\ran}{\operatorname{ran}}
\newcommand{\dom}{\operatorname{dom}}
\newcommand{\tc}{\operatorname{tc}}
\newcommand{\Sat}{\operatorname{Sat}}
\newcommand{\Pow}{\mathcal P}
\newcommand{\forces}{\Vdash}
\newcommand{\PP}{\mathbb P}
\newcommand{\QQ}{\mathbb Q}
\newcommand{\seq}[1]{\langle #1\rangle}
\newcommand{\tail}{\operatorname{tail}}
\newcommand{\Cof}{\operatorname{Cof}}
\newtheorem{remark}[definition]{Remark}
\newtheorem{question}[definition]{Question}
\newtheorem*{maintheorem}{Main theorem}
\newtheorem*{widththeorem}{Definable-family theorem}
\begin{document}
\thispagestyle{empty}
{\sffamily\bfseries\small\color{Olive}SET THEORY\quad /\quad RESEARCH CONTINUATION}

\vspace{13pt}
{\fontsize{28}{31}\selectfont\bfseries\color{Forest}Large cardinals at the\\ inconsistency frontier\par}
\vspace{10pt}
{\fontsize{14}{18}\selectfont A Prikry consistency refinement and\\ local definability obstructions\par}
\vspace{14pt}
{\color{Muted}Prepared for Vladimir\hfill 18 September 2026\par}
\vspace{3pt}
{\color{Sage}\hrule height .6pt}
\vspace{12pt}

\textbf{Principal result developed here.} Let $j:V_\Lambda\to V_\Lambda$ be an
$\Iwf$-embedding, let $\kappa=\crit(j)$, and let $U$ be its derived normal
ultrafilter. Prikry forcing with $U$ makes $\kappa$ \emph{$\kappa$-cover exacting}.
Consequently,
\[
 \boxed{\quad
 \Con(\ZFC+\Iwf)\ \Longrightarrow\
 \Con\bigl(\ZFC+\exists\lambda\,\CEx_\lambda(\lambda)\bigr).
 \quad}
\]
This replaces the supplied report's explicit $\Itwo$ upper bound by a weaker
rank-into-rank assumption. The construction combines the well-founded-part
localization used for ordinary exactingness with a finite-change family of
Prikry-name evaluations. A separate tree argument puts the required local
embedding \emph{inside} the constructed model.

\begin{assessment}{Positive and negative conclusions}
\textbf{Relative consistency.} Sections~\ref{sec:internal}--\ref{sec:forcing}
give the construction and its internalization argument, including the
uniform cardinal bound on the covering family.

\textbf{Conditional inconsistency.} If $\lambda$ is exacting, no nonempty
$\OD_{V_\lambda}$ family of fewer than $\lambda$ short cofinal maps into
$\lambda$ exists. In fact, one coordinate projection of any such definable
family has order type $\lambda$. Individual maps need not be definable.

\textbf{Scope.} These results do not settle cover exactingness above a
strongly compact cardinal, ordinary exactingness above an extendible, or the
unrestricted cardinal-preserving embedding problem.
\end{assessment}

\textbf{Provenance and status.} This is a research draft with detailed proofs,
not a claim of an independently refereed discovery. The iteration facts are
from Andretta--Dimonte and Aguilera--Bagaria--Goldberg--L\"ucke. Blue--Goldberg's
2026 notes already anticipate cover versions of related constructions.
The precise refinement below is established here as a continuation of the
supplied report; its priority in the literature is not certified.
No proof-assistant verification is claimed.

\clearpage
\tableofcontents
\vspace{8pt}
\begin{assessment}{Reading route}
The main forcing theorem is Theorem~\ref{thm:forcing}; the two reusable
arguments are the internalization lemma, Lemma~\ref{lem:internal}, and the
finite-change cloud lemma, Lemma~\ref{lem:cloud}. The local negative result
is Theorem~\ref{thm:width}. Section~\ref{sec:audit} records the proof's
sensitive interfaces and the questions it leaves open.
\end{assessment}

\clearpage
\section{What this continuation adds}\label{sec:scope}

The supplied unified report separates an axiom from an additional closure,
definability, or stationarity requirement that might make it contradictory.
It also singles out cover exactingness as a useful boundary: its existence
has a consistency construction, while placing it above an extendible
cardinal is ruled out in the cited lecture notes. We keep that organization,
but pursue a positive construction and a more flexible local obstruction
rather than repeat the candidate survey.\cite{unified,cover}

The starting comparison in the supplied report was
$\Con(\ZFC+\Itwo)\Rightarrow\Con(\ZFC+\exists\lambda\,\CEx_\lambda(\lambda))$.
The ordinary-exacting construction in Aguilera--Bagaria--Goldberg--L\"ucke
uses only $\Iwf$.\cite[Theorem~5.4]{abgl}
The mathematical task is to retain a \emph{uniformly small predicate}
containing an arbitrary object, not just an unexpanded exacting witness.

\begin{maintheorem}
Suppose $j:V_\Lambda\to V_\Lambda$ is an $\Iwf$-embedding with critical
point $\kappa$, and
\[
 U=\{A\subseteq\kappa:\kappa\in j(A)\}.
\]
If $G$ is generic for Prikry forcing $\PP_U$, then
\[
 V[G]\models\CEx_\kappa(\kappa).
\]
In particular, the existence of an $\Iwf$-embedding implies the existence
of a countable transitive model of $\ZFC$ with a $\lambda$-cover-exacting
cardinal.
\end{maintheorem}

The last assertion is stronger than an informal appeal to a generic over
the universe. Its proof passes to a countable elementary submodel after
establishing the forcing theorem; see Corollary~\ref{cor:transitive}.
No converse consistency reduction is asserted.

The second result strengthens the original report's single-definable-map
obstruction. Write $\OD_{V_\lambda}$ for definability from ordinal parameters
and finitely many parameters belonging to $V_\lambda$.

\begin{widththeorem}
Let $\lambda$ be exacting and let $0<\rho<\lambda$. Suppose
$F\in\OD_{V_\lambda}$ is a nonempty family of cofinal maps
$\rho\to\lambda$. Then for some $\xi<\rho$,
\[
 \ot\{f(\xi):f\in F\}=\lambda.
\]
Consequently $|F|\geq\lambda$ in the ambient universe. If every member
of $F$ is increasing, all coordinate projections from some coordinate
onward have order type $\lambda$.
\end{widththeorem}

The hypothesis is that the \emph{family} is definable. It is not assumed
that a member of the family belongs to HOD, that the family has a definable
well-order, or that a definable choice function selects one of its members.
This distinction is what makes the obstruction stronger than simply
asking for a cofinal map in $\HOD_{V_\lambda}$.

\subsection{A precise novelty ledger}

\begin{center}
\small
\renewcommand{\arraystretch}{1.18}
\begin{tabularx}{\textwidth}{@{}P{.29\textwidth}Y@{}}
\toprule
\textbf{Ingredient or conclusion}&\textbf{Status in this continuation}\\\midrule
Iteration and well-founded-part facts&Imported mathematical inputs, with the exact source locators recorded below.\\
The $\Itwo$ cover construction&Already in Blue--Goldberg's notes; not presented as a discovery here.\\
The $\Iwf$ cover refinement&An explicit proof developed here by combining and extending those arguments. The notes' broader anticipation of cover variants prevents a justified priority claim.\\
The definable-family obstruction&A proved consequence of the established HOD-regularity theorem, with a quantitative projection bound. It is not a new refutation of bare exactingness.\\
The smaller-strongly-compact case&Unresolved by these arguments. No preservation of such a cardinal below the new exacting cardinal is proved.\\
\bottomrule
\end{tabularx}
\end{center}

A distinction between a new deduction in this document and a new theorem
in the historical literature is essential. The latter requires a more
exhaustive priority investigation and specialist scrutiny than a targeted
source check can provide. The mathematical claims below are stated with
their hypotheses and proofs, rather than with a claim of priority.

\section{Definitions and the iteration interface}\label{sec:definitions}

\subsection{Cover exactingness and its quantifiers}

\begin{definition}\label{def:relative}
For $\alpha>\lambda$ and $Y\subseteq V_\alpha$, let
$\RelEx(\lambda,\alpha,Y)$ assert that there are $X$ and an elementary map
\[
 k:(X,\in,Y\cap X)\longrightarrow(V_\alpha,\in,Y)
\]
with
\[
 (X,\in,Y\cap X)\prec(V_\alpha,\in,Y),\qquad
 V_\lambda\cup\{\lambda\}\subseteq X,\qquad
 k(\lambda)=\lambda,\qquad k\restr\lambda\ne\id.
\]
A cardinal $\lambda$ is \emph{$\gamma$-cover exacting} if
\begin{equation}\label{eq:cover}
 \forall\alpha>\lambda\ \forall y\in V_\alpha\ \exists Y\subseteq V_\alpha
 \quad
 \bigl[y\in Y\ \wedge\ |Y|\leq\gamma\ \wedge\
       \RelEx(\lambda,\alpha,Y)\bigr].
\end{equation}
\end{definition}

This is the bounded-predicate definition used in the supplied report and
in Blue--Goldberg's notes.\cite[Definitions~2.6 and~3.1]{cover}
The object $y$ is a \emph{member} of $Y$. It need not belong to $X$, and
the witness need not fix $y$. The bound $\gamma$ is independent of
$\alpha$ and $y$. The substructure $X$ need not be transitive.
Every occurrence of this definition in a model $W$ uses $V_\alpha^W$,
cardinality in $W$, and witnesses belonging to $W$.

Ordinary exactingness, $\Ex(\lambda)$, requires the same witness after
omitting the predicate $Y$, for every cardinal height $\alpha>\lambda$.
Every cover-exacting cardinal is therefore exacting.

The language is finite, so satisfaction for these set structures is
available in ordinary set theory. Thus \eqref{eq:cover} is a first-order
set-theoretic assertion, not a proper-class embedding axiom.

\subsection{The weak iterability hypothesis}

Fix $j:V_\Lambda\to V_\Lambda$ with critical sequence
\[
 \kappa_0=\crit(j),\qquad \kappa_{n+1}=j(\kappa_n),\qquad
 \Lambda=\sup_{n<\omega}\kappa_n.
\]
All uses of $\Ithree$ in the construction have this normalization. The
canonical extension on a subset $A\subseteq V_\Lambda$ is
\[
 j^+(A)=\bigcup_{\beta<\Lambda}j(A\cap V_\beta).
\]
Let $j_{0,1}=j$, let
$j_{n+1,n+2}=j^+(j_{n,n+1})$, and form the commuting finite iteration
system $\seq{j_{m,n}:m\leq n<\omega}$. Its direct limit is denoted
$M^j_\omega$, with maps $j_{n,\omega}$. The limit may be ill-founded.
Its well-founded part, identified with its transitive collapse, is
$W^j_\omega$.

\begin{definition}\label{def:wf}
The embedding $j$ is an $\Iwf$-embedding if
\begin{equation}\label{eq:wf}
 j_{0,\omega}[\kappa_0^+]\subseteq W^j_\omega.
\end{equation}
\end{definition}

This is Definition~5.2 of Aguilera--Bagaria--Goldberg--L\"ucke with
$n=0$.\cite{abgl} It does \emph{not} require the whole direct limit to
be well-founded. In particular, it is not the same hypothesis as
$\Iwfone$, which requires full well-foundedness of the first
$\omega$-iterate.

\subsection{The imported interface, rather than an assumed extension}

We use the following iteration facts. The finite-iteration construction
and amenable-predicate elementarity are treated by
Andretta--Dimonte; the fixed-code and well-founded-part formulations below
are recorded in Aguilera--Bagaria--Goldberg--L\"ucke.\cite[Section~3,
 especially Lemma~3.3]{ad}\cite[Proposition~5.1 and proof of Theorem~5.4]{abgl}

\begin{proposition}[Iteration interface]\label{prop:interface}
With the preceding notation:
\begin{enumerate}
\item $V_\Lambda\cup\{\Lambda\}\subseteq W^j_\omega$,
$j_{n,\omega}\restr V_{\kappa_n}=\id$, and
$j_{n,\omega}(\kappa_n)=\Lambda$.
\item If $A\in V_{\kappa_0+1}$, then
$j_{0,\omega}(A)\in V_{\Lambda+1}\cap W^j_\omega$ and
$j^+(j_{0,\omega}(A))=j_{0,\omega}(A)$.
\item If $E\subseteq\Lambda\times\Lambda$ and $j^+(E)=E$, then
$j\restr\Lambda:(\Lambda,E)\to(\Lambda,E)$ is elementary.
\item Put $U_0=\{A\subseteq\kappa_0:\kappa_0\in j(A)\}$ and
$U=j_{0,\omega}(U_0)$. For $A\in\Pow(\Lambda)\cap W^j_\omega$,
\begin{equation}\label{eq:tailmeasure}
 A\in U\quad\Longleftrightarrow\quad
       \exists m\ \forall n\geq m\ (\kappa_n\in A).
\end{equation}
\end{enumerate}
\end{proposition}

Clause (3) concerns the \emph{expanded} structure, not an assertion that
$j^+$ is an $\Ione$-embedding. One applies amenable-predicate elementarity
to $(V_\Lambda,E)$ and then interprets the ordinal domain. Every bounded
piece of $E$ belongs to $V_\Lambda$, and $\cf(\Lambda)=\omega$ supplies
the relevant hypothesis of the cited lemma.

For clause (4), a subset represented at a finite stage has its eventual
membership along the critical sequence determined by that stage's derived
normal measure. If it is measure one, all sufficiently late critical
points belong to it; if it is measure zero, the complementary set is
measure one. This also explains why \eqref{eq:tailmeasure} is an
\emph{external} description of an internal ultrafilter. It does not
make $\Lambda$ measurable in the ambient universe, where
$\cf(\Lambda)=\omega$.

These are the only rank-into-rank iteration inputs needed below. The new
internalization argument does not assume an extension of $j$ to an
arbitrary higher rank or to the universe.

\section{Internalizing an external local witness}\label{sec:internal}

The main obstacle is not the existence of an external map. A forcing
construction has to produce a model that \emph{contains} the embedding
witnesses in Definition~\ref{def:relative}. The following lemma handles
this transfer, including high-rank predicates.

\begin{lemma}[Enumerated-domain internalization]\label{lem:internal}
Work in an ambient universe satisfying $\ZFC$. Let $W$ be a transitive
set model of a sufficiently strong finite fragment of $\ZFC$. Suppose
there are an external elementary map $e:W\to W$, an infinite cardinal
$\lambda$ of $W$, and $c=\seq{\nu_n:n<\omega}\in W$ such that
\begin{enumerate}
\item $c$ is strictly increasing and cofinal in $\lambda$,
$e(\lambda)=\lambda$, and $e(\nu_n)=\nu_{n+1}$;
\item $\omega<\nu_0$, $V_\lambda^{\mathrm{ambient}}\subseteq W$, and
$W\models |V_\lambda|=\lambda$;
\item $\alpha>\lambda$ and $Y\subseteq V_\alpha^W$ belong to $W$,
with $e(\alpha)=\alpha$ and $e(Y)=Y$.
\end{enumerate}
Then
\[
 W\models\RelEx(\lambda,\alpha,Y).
\]
It is not assumed that $e\in W$ or that $e\restr V_\lambda\in W$.
\end{lemma}

\begin{proof}
All the structures, functions, and trees that are said to be chosen
\emph{in $W$} below are sets of $W$. Write
\[
 \mathcal A=(V_\alpha^W,\in,Y).
\]
Inside $W$, choose an elementary substructure $X\prec\mathcal A$ of
cardinality $\lambda$ containing $V_\lambda^W\cup\{\lambda\}$.
The downward L\"owenheim--Skolem construction applies to the finite
language with its unary predicate. Set $Z=e(X)$, with the induced
expanded structure. Since $e(\mathcal A)=\mathcal A$,
\[
 W\models Z\prec\mathcal A.
\]
In $W$, take a bijection $f:\lambda\to X$ with
$f(0)=\lambda$ and $f(1)=\nu_0$, and put $g=e(f):\lambda\to Z$.
Then $g(0)=\lambda$ and $g(1)=\nu_1$.

For $n<\omega$, define in $W$
\[
 X_n=f``\nu_n,\qquad Z_n=g``\nu_{n+1}.
\]
Let $T$ be the following tree in $W$. A node at level $n$ is a finite
list $(p_0,\ldots,p_n)$ such that, for every $m\leq n$,
\begin{enumerate}
\item $p_m:X_m\to Z_m$ is a partial elementary map from the full
structure on $X$ to the full structure on $Z$;
\item $p_m(\lambda)=\lambda$ and $p_m(\nu_0)=\nu_1$;
\item $p_{m+1}\restr X_m=p_m$ when $m<n$.
\end{enumerate}
The order is end extension of finite lists. Here ``partial elementary''
means agreement on every first-order formula whose finitely many
parameters lie in the map's domain. It does not mean elementarity between
$X_m$ and $Z_m$ as separate substructures. Satisfaction for $X$ and $Z$
is a set in $W$, so this definition gives a set tree in $W$.

For each $n$, consider the external restriction
\[
 q_n=e\restr\nu_n.
\]
It maps $\nu_n$ into $\nu_{n+1}$, and its graph has rank strictly below
$\lambda$. Hence $q_n\in V_\lambda^{\mathrm{ambient}}\subseteq W$.
Using $f,g,q_n$ inside $W$, form
\begin{equation}\label{eq:partialmaps}
 p_n(f(\xi))=g(q_n(\xi))\qquad(\xi<\nu_n).
\end{equation}
The elementary evaluation identity gives, externally,
\[
 g(e(\xi))=e(f)(e(\xi))=e(f(\xi)),
\]
so $p_n=e\restr X_n$. Thus each $p_n$ is partial elementary and has
the required two prescribed values. Satisfaction for the transitive
model's set structures is absolute, so $W$ recognizes these properties.
The finite lists $(p_0,\ldots,p_n)$ belong to $T$. They form an external
infinite branch.

If $W$ had no infinite branch through $T$, its well-founded recursion
and dependent choice would give a rank function on the strict-extension
relation of $T$ with values in its ordinals. Since $W$ is transitive,
that would be an actual ordinal-valued rank function in the ambient
universe. It would strictly decrease along the displayed external branch,
which is impossible. Therefore a branch $B$ belongs to $W$.

The union of the maps along $B$ is an elementary map $k:X\to Z$ in $W$.
Its domain is all of $X$, because $\sup_n\nu_n=\lambda$ and $f$ is
onto. Every finite tuple is contained in one $X_n$, which verifies full
elementarity of the union. Moreover,
\[
 k(\lambda)=\lambda,\qquad k(\nu_0)=\nu_1\ne\nu_0.
\]
Composing with the elementary inclusion $Z\prec\mathcal A$ proves
$\RelEx(\lambda,\alpha,Y)$ in $W$.
\end{proof}

\subsection{Why the argument is genuinely internal}

Two tempting shortcuts would be wrong. The external map $e\restr X$
is not declared to belong to $W$. Instead, each bounded restriction is
\emph{recoded by a small ordinal map} $q_n$ that does belong to $W$.
Nor is the external branch simply imported into $W$: well-foundedness
absoluteness supplies some internal branch, possibly a different one.

The finite-list definition of $T$ is also deliberate. A map
$X_n\to Z_n$ need not have its restriction to $X_m$ land in $Z_m$ for
$m<n$. Requiring a coherent finite list, with a separate range bound at
each level, makes the tree closed under initial segments and avoids this
otherwise easy oversight.

The required fragment of set theory can be fixed uniformly before the
application. It includes the set-satisfaction construction, the elementary
hull construction, the indicated instances of Replacement and Separation,
and the theorem that a well-founded set relation has an ordinal rank.
Appendix~\ref{app:fragment} explains how this fragment is secured by
reflection. No assumption that an arbitrary rank segment satisfies all
of $\ZFC$ is made.

\section{Finite-change clouds of Prikry names}\label{sec:clouds}

Let $N$ be transitive, let $D\in N$ be a normal measure on $\lambda$
\emph{as computed in $N$}, and let $\QQ=\PP_D^N$. For an increasing
cofinal sequence $h:\omega\to\lambda$, define
\[
 G_h=\{(s,A)\in\QQ:s\text{ is an initial segment of }h
                  \text{ and }h[\omega\setminus\dom(s)]\subseteq A\}.
\]
The forcing order is the usual one: extend the finite increasing stem,
choose its new points from the measure-one set, and shrink that set.
The parameter $\QQ$ denotes the ground partial order. It is not a newly
computed forcing from all subsets of $\lambda$ in the extension.

For increasing sequences $c,h$, write
\begin{equation}\label{eq:finitechange}
 h\sim c\quad\Longleftrightarrow\quad
       |\ran(h)\mathbin{\triangle}\ran(c)|<\omega.
\end{equation}
This is equality modulo finite change of the \emph{ranges}, not
coordinatewise eventual equality of sequences. In particular,
$c^-=\seq{c(n+1):n<\omega}$ satisfies $c^-\sim c$.

\Needspace{185pt}
\begin{lemma}[Finite-change cloud]\label{lem:cloud}
Suppose $c$ is $\PP_D^N$-generic over $N$ and $W=N[G_c]$. Let
$\tau\in N$ be a $\QQ$-name, and suppose
\[
 N\models\mathbf 1_{\QQ}\forces\tau\in V_{\check\alpha}.
\]
Define in $W$
\begin{equation}\label{eq:cloud}
 Y_{\tau,c}=
 \{\tau^{G_h}:h\in{}^\omega\lambda\text{ is strictly increasing and }h\sim c\}.
\end{equation}
Then $Y_{\tau,c}\in W$, $\tau^{G_c}\in Y_{\tau,c}\subseteq V_\alpha^W$,
and $W\models |Y_{\tau,c}|\leq\lambda$.

If an external elementary $e:W\to W$ fixes $\lambda,\QQ,\tau,\alpha$
and satisfies $e(c)=c^-$, then $e(Y_{\tau,c})=Y_{\tau,c}$.
\end{lemma}

\begin{proof}
Every $h\sim c$ is eventually in every member of $D$, because $c$ has
that property. The Mathias criterion therefore makes $h$ generic over
$N$. This criterion and the elementary properties of Prikry forcing are
used in their model-relative forms.\cite[proof of Theorem~5.4]{abgl}

Moreover, $N[G_h]=N[G_c]$. Indeed, $h\in W$ gives one inclusion.
For the other, the two finite difference sets are finite sets of
ordinals below $\lambda$, hence are elements of $N$. Using these sets,
$c$ is recovered from $h$ by taking the increasing enumeration of its
modified range. Thus $c\in N[G_h]$.
It follows from the forcing theorem that every value in \eqref{eq:cloud}
belongs to the same $V_\alpha^W$.

The definition of $Y_{\tau,c}$ inside $W$ uses only $\tau,\QQ,c$ and
$\lambda$. It does \emph{not} quantify over $N$ or require that $N$ be
an element of $W$. A forcing-name evaluation is a well-founded recursive
operation on the name and the set $G_h$. Replacement in $W$ therefore
forms the displayed family.

An $h\sim c$ is determined by two finite subsets of $\lambda$: points
deleted from $\ran(c)$ and points added to it. There are at most
$\lambda$ such pairs in $W$, since $\lambda$ is an infinite cardinal.
Consequently there are at most $\lambda$ possible $h$ and at most
$\lambda$ evaluations. The choice $h=c$ includes the required object.

Finally, elementarity applied to the internal definition gives
\[
 e(Y_{\tau,c})=Y_{e(\tau),e(c)}=Y_{\tau,c^-}.
\]
The relation $\sim$ is an equivalence relation, and $c^-\sim c$, so the
defining families of sequences are identical. Hence
$Y_{\tau,c^-}=Y_{\tau,c}$.
\end{proof}

\begin{remark}
The value $y=\tau^{G_c}$ need not be fixed. In fact,
$e(y)=\tau^{G_{c^-}}$ may be a different member of $Y_{\tau,c}$.
Cover exactingness is designed to tolerate precisely this motion inside
an invariant small family. Replacing that family by the singleton
$\{y\}$ would invalidate the argument.
\end{remark}

The finite-change device is related to the name families in
Blue--Goldberg's construction.\cite[proof of Theorem~3.5]{cover}
Using just the explicitly defined finite-change class has a useful
advantage here: we do not need a converse theorem classifying
\emph{all} Prikry generics that yield the same extension.

\section{The Prikry consistency refinement}\label{sec:forcing}

\subsection{A localized shift, with fixed name parameters}

For the rest of this section, $j:V_\Lambda\to V_\Lambda$ is an $\Iwf$
embedding, $\kappa=\kappa_0$, and $U_0$ is the derived normal measure on
$\kappa$. The following localization is the bridge from the iteration
interface to the two lemmas just proved.

\Needspace{230pt}
\begin{lemma}[Coded localization and parameter fixedness]\label{lem:localize}
Let $N_0$ be a transitive set of cardinality $\kappa$ satisfying a fixed,
sufficiently strong finite fragment of $\ZFC$, with
\[
 V_\kappa\cup\{\kappa,\bar U_0\}\subseteq N_0,
 \qquad \bar U_0=U_0\cap N_0,
 \qquad \Ord\cap N_0<\kappa^+.
\]
Suppose $N_0$ regards $\bar U_0$ as a normal measure on $\kappa$.
Then there are a transitive set $N$, an elementary map $i:N\to N$,
and the corresponding map $J=j_{0,\omega}$ on $N_0$-parameters such that
\begin{enumerate}
\item $N=J(N_0)$ is well-founded,
$V_\Lambda\cup\{\Lambda\}\subseteq N$, and
$i\restr\Lambda=j\restr\Lambda$;
\item $D=J(\bar U_0)\in N$ is a normal measure on $\Lambda$ in $N$;
\item $i(J(a))=J(a)$ for every $a\in N_0$;
\item $c=\seq{\kappa_n:n<\omega}$ is Prikry-generic over $N$ for
$D$, and $i$ lifts to an elementary map
\[
 e:W=N[G_c]\longrightarrow W
 \quad\text{with}\quad e(c)=c^-,\qquad e\restr N=i.
\]
\end{enumerate}
In particular, $W$ contains the full ambient $V_\Lambda$ and regards
$\Lambda$ as an infinite cardinal with $|V_\Lambda|=\Lambda$.
\end{lemma}

\begin{proof}
We spell out the parameter issue in the well-founded coding construction
used in the proof of \cite[Theorem~5.4]{abgl}. It matters because the
Prikry name selected later need not be ordinal definable.

Choose a bijection $b_0:\kappa\to N_0$ satisfying
\begin{equation}\label{eq:ordinalslots}
 b_0(0)=\kappa,\qquad b_0(1)=\bar U_0,\qquad
 b_0(\omega\cdot(\beta+1))=\beta\quad(\beta<\kappa).
\end{equation}
The reserved ordinal slots are disjoint from $0,1$; both their complement
and the remaining values have cardinality $\kappa$. This permits completion
to a bijection. The successor in $\beta+1$ avoids a collision between the
slot for the ordinal $0$ and the slot for $\kappa$.
Let
\[
 \xi\,E_0\,\eta\quad\Longleftrightarrow\quad b_0(\xi)\in b_0(\eta).
\]
The sets $N_0,b_0,E_0$ have rank below $\kappa^+$ and hence belong to
$V_\Lambda$. Define, in the direct limit,
\[
 E=J(E_0),\quad N=J(N_0),\quad b=J(b_0),\quad D=J(\bar U_0).
\]
If $\delta=\Ord\cap N_0$, the hypothesis \eqref{eq:wf} puts
$J(\delta)$ in the well-founded part. The internal rank function of the
transitive set $N$ is bounded by this genuine ordinal. Thus $N$ and its
membership relation are well-founded, and their collapse is legitimate.
We identify $N$ with that collapse. The iteration interface gives
$V_\Lambda\cup\{\Lambda\}\subseteq N$.

Now $b:(\Lambda,E)\cong(N,\in)$, and $j^+(E)=E$ by
Proposition~\ref{prop:interface}. Consequently
\[
 i=b\circ(j\restr\Lambda)\circ b^{-1}:N\to N
\]
is elementary. Applying $J$ to \eqref{eq:ordinalslots} and using the
absoluteness of ordinal arithmetic gives
\[
 b(\omega\cdot(\beta+1))=\beta\qquad(\beta<\Lambda).
\]
Therefore $i(\beta)=j(\beta)$ for $\beta<\Lambda$.

Here is the additional parameter observation. Given any $a\in N_0$,
choose $\xi<\kappa$ with $b_0(\xi)=a$. Since $J(\xi)=j(\xi)=\xi$,
\begin{equation}\label{eq:fixedrange}
 J(a)=b(\xi),\qquad i(J(a))=b(j(\xi))=b(\xi)=J(a).
\end{equation}
Thus the shift fixes \emph{every image parameter from $N_0$}, not merely
$\Lambda$ and $D$. This does not make $i$ the identity on $N$:
$J``N_0$ need not be all of $N$.

Elementarity gives that $D$ is a normal measure on $\Lambda$ in $N$,
with $D=N\cap J(U_0)$. By \eqref{eq:tailmeasure}, every member of $D$
contains a tail of $c$. The Mathias criterion makes $G_c$ generic over
$N$. The tail $c^-$ is also generic and generates the same extension.
Since $i(\QQ)=\QQ$ for $\QQ=\PP_D^N$ and $i``G_c\subseteq G_{c^-}$,
the usual name interpretation defines the lift
\begin{equation}\label{eq:lift}
 e(\sigma^{G_c})=i(\sigma)^{G_{c^-}}.
\end{equation}
The forcing theorem proves that this is well-defined and elementary.
The canonical generic-sequence name is fixed by $i$, so $e(c)=c^-$.

Prikry forcing preserves $\Lambda$ as a cardinal in $N[G_c]$ and adds
no bounded subsets of it. Also $V_\Lambda^{\mathrm{ambient}}\subseteq N
\subseteq W\subseteq V$, so the rank $V_\Lambda$ is the same in $W$
and the ambient universe. Each bounded rank has cardinality below
$\Lambda$ in $W$: choose a larger $\kappa_n$ and use a bounding
bijection of rank below $\Lambda$, which already belongs to $W$.
Since $c\in W$ is cofinal in $\Lambda$, Choice in $W$ gives
$|V_\Lambda|^W=\Lambda$.
\end{proof}

\subsection{Homogeneity and one named counterexample}

For later use define the first-order formula
\begin{equation}\label{eq:bad}
\begin{split}
 \Bad(\kappa,\alpha,y)\quad\Longleftrightarrow\quad
 &\alpha>\kappa\ \wedge\ y\in V_\alpha\ \wedge\
 \\[-2pt]
 &\forall Y\subseteq V_\alpha\,
 \bigl[(y\in Y\wedge |Y|\leq\kappa)
               \ \Longrightarrow\ \neg\RelEx(\kappa,\alpha,Y)\bigr].
\end{split}
\end{equation}
For a cardinal $\kappa$, $\neg\CEx_\kappa(\kappa)$ is equivalent to
$\exists\alpha\exists y\,\Bad(\kappa,\alpha,y)$.
Prikry forcing preserves $\kappa$ as a cardinal, so this formulation
applies in the extensions under consideration.

We need only the following familiar consequence of Prikry homogeneity:
a sentence about the extension using ground-model parameters, but not a
name for the chosen generic, is decided by the top condition.
For completeness, below two Prikry conditions one can shrink the
measure-one parts to a common set above both stems. Their cones are
isomorphic by replacing the fixed stem and retaining the appended tail.
The isomorphism preserves canonical names for ground-model parameters.
Opposite decisions of such a sentence would therefore be impossible.
This cone-homogeneity argument is enough; no particular global
permutation of the entire forcing is needed.

If the top condition forces failure of cover exactingness, the least
height $\alpha$ of failure is definable in the extension from $\kappa$.
Some condition decides its ordinal value. Applying the preceding
homogeneity fact to that value shows that a single ground-model ordinal
$\alpha$ is forced to be the least bad height by the top condition.
The maximum principle then supplies a name $\tau$ with
\begin{equation}\label{eq:badforced}
 \mathbf 1_{\PP_{U_0}}\forces\Bad(\check\kappa,\check\alpha,\tau).
\end{equation}
The name $\tau$ itself need not be invariant under forcing isomorphisms.
Its later invariance under the \emph{localized shift} follows instead
from \eqref{eq:fixedrange}.

\subsection{The complete forcing argument}

\begin{theorem}\label{thm:forcing}
If $j:V_\Lambda\to V_\Lambda$ is an $\Iwf$-embedding, $\kappa=\crit(j)$,
and $U_0$ is the derived normal measure, then
\begin{equation}\label{eq:forcingtheorem}
 \mathbf 1_{\PP_{U_0}}\forces\CEx_{\check\kappa}(\check\kappa).
\end{equation}
\end{theorem}

\begin{proof}
Suppose otherwise. Homogeneity and the preceding discussion supply a
height $\alpha$ and a name $\tau$ satisfying \eqref{eq:badforced}.
Choose a reflecting rank $V_\zeta$ above $\Lambda,\alpha$ and the ranks
of the names and forcing data. Require reflection of
\eqref{eq:badforced}, of the fixed finite fragment needed for the
preceding lemmas, and of the forcing facts used in them.
Inside the ambient universe choose
\[
 H\prec V_\zeta,\qquad |H|=\kappa,\qquad
 V_\kappa\cup\{\kappa,U_0,\alpha,\tau\}\subseteq H.
\]
Let $\pi:H\cong N_0$ be the transitive collapse. It fixes $V_\kappa$
and $\kappa$. Put
\[
 \bar U_0=\pi(U_0),\qquad a_0=\pi(\alpha),\qquad t_0=\pi(\tau).
\]
For subsets of $\kappa$ that belong to $H$, the collapse is the identity.
It follows that $\bar U_0=N_0\cap U_0$. The ordinal height of $N_0$
has cardinality at most $\kappa$ and is below $\kappa^+$.
Thus Lemma~\ref{lem:localize} applies.

Let $J=j_{0,\omega}$ and use the resulting $N,i,D,c,W,e$.
Set
\[
 a=J(a_0),\qquad t=J(t_0),\qquad \QQ=\PP_D^N.
\]
The parameter fixedness clause gives
\begin{equation}\label{eq:fixedparameters}
 e(a)=a,\qquad e(t)=t,\qquad e(\QQ)=\QQ,\qquad e(\Lambda)=\Lambda.
\end{equation}
Elementarity of the collapse and of the direct-limit map transfers
\eqref{eq:badforced} to
\[
 N\models \mathbf 1_{\QQ}\forces\Bad(\check\Lambda,\check a,t).
\]
This is a statement about the \emph{internal} forcing relation of $N$.
Since $G_c$ is $N$-generic, its forcing theorem gives
\begin{equation}\label{eq:badinw}
 W\models\Bad(\Lambda,a,t^{G_c}).
\end{equation}

Apply Lemma~\ref{lem:cloud} to $t$ and $c$. In $W$ it produces the set
\[
 Y=\{t^{G_h}:h\sim c\},
 \qquad t^{G_c}\in Y\subseteq V_a^W,\qquad |Y|^W\leq\Lambda.
\]
Here and below the sequences in the cloud are strictly increasing.
Equations~\eqref{eq:fixedparameters} and $e(c)=c^-$ imply $e(Y)=Y$.
The hypotheses of Lemma~\ref{lem:internal} now hold for
\[
 (W,e,\lambda,c,\alpha,Y)=(W,e,\Lambda,c,a,Y).
\]
In particular, $e(\kappa_n)=\kappa_{n+1}$, the full ambient
$V_\Lambda$ is in $W$, and $W$ regards its size as $\Lambda$.
Therefore
\[
 W\models\RelEx(\Lambda,a,Y).
\]
This contradicts \eqref{eq:badinw}, because the same $Y$ belongs to $W$,
contains $t^{G_c}$, and has size at most $\Lambda$ there.
The contradiction proves \eqref{eq:forcingtheorem}.
\end{proof}

\begin{assessment}{Where the strengthening occurs}
The localized iteration alone provides an external self-embedding of
$W$. The cloud lemma turns an arbitrary name value into a small invariant
predicate, and the enumerated-domain tree turns preservation of that
predicate into an internal witness at the required height. Neither
``the embedding exists externally'' nor ``the name has a fixed code''
would suffice by itself.
\end{assessment}

\section{Consistency consequences and exact limits}\label{sec:consequences}

\subsection{A set-model conclusion}

\begin{corollary}\label{cor:transitive}
The existence of an $\Iwf$-embedding implies the existence of a countable
transitive model of
\[
 \ZFC+\exists\lambda\,\CEx_\lambda(\lambda).
\]
Consequently,
\[
 \Con(\ZFC+\Iwf)\ \Longrightarrow\
 \Con\bigl(\ZFC+\exists\lambda\,\CEx_\lambda(\lambda)\bigr).
\]
\end{corollary}

\begin{proof}
For a normalized rank-into-rank embedding, $V_\Lambda$ is a model of
$\ZFC$. We use the rank-model passage from
\cite[proof of Corollary~5.5]{abgl}, checking the additional cover
witnesses explicitly below. The forcing $\PP_{U_0}$ and all its subsets have rank below
$\Lambda$. A generic over $V_\Lambda$ is therefore generic over $V$,
and
\[
 V_\Lambda[G]=V[G]_\Lambda.
\]
By Theorem~\ref{thm:forcing}, the latter rank model satisfies
$\CEx_\kappa(\kappa)$. To check the downward passage explicitly, a
witness at height $\alpha<\Lambda$ consists of $Y\subseteq V_\alpha$,
$X\subseteq V_\alpha$, and the graph of a map between subsets of
$V_\alpha$. Their ranks are at most $\alpha$ plus a fixed finite amount,
which is below $\Lambda$. The same is true of a bijection certifying
$|Y|\leq\kappa$. Set-structure satisfaction is absolute. Thus the
witnesses can be retained in $V[G]_\Lambda$.

It follows that $V_\Lambda$ satisfies the internal forcing assertion
that its top condition forces $\CEx_\kappa(\kappa)$. Choose a countable
$A\prec V_\Lambda$ containing $U_0,\kappa$ and collapse it to a
transitive model $M$. This forcing assertion transfers to $M$.
A generic filter for its collapsed forcing exists in the ambient universe
by enumerating the countably many dense sets of $M$. Its extension is a
countable transitive model of the displayed theory.

This is an ordinary relative-consistency argument formalizable over the
stated theories. It does not require an external assumption that an
arbitrary consistent theory already has a transitive model.
\end{proof}

\subsection{The revised strength interval}

Use $A\succeq_{\rm con}B$ for the indicated relative-consistency
reduction, and use strict comparison only where it is supplied by the
established iteration hierarchy. The result yields
\begin{equation}\label{eq:strengthchain}
\begin{gathered}
 \Itwo\ \succ_{\rm con}\ \Iwf
 \ \succeq_{\rm con}\
       \exists\lambda\,\CEx_\lambda(\lambda),\\
 \exists\lambda\,\CEx_\lambda(\lambda)
 \ \succeq_{\rm con}\ \exists\lambda\,\Ex(\lambda)
 \ \succ_{\rm con}\ \Ithree.
\end{gathered}
\end{equation}
The strict comparisons and the ordinary-exacting lower bound are those
of the existing hierarchy.\cite[Theorem~5.10 and Corollary~5.11]{abgl}
The two middle comparisons are \emph{not} claimed to be strict, and
no equiconsistency between exactingness and cover exactingness follows.

In particular, the least-budget cover principle is placed strictly
between $\Ithree$ and $\Itwo$ in this consistency-strength sense.
The additional uniform cover requirement has not forced the upper bound
back up to $\Itwo$.

\subsection{The cardinal-ordering issue survives unchanged}

The argument does not preserve a specified strongly compact cardinal
below $\kappa$ through the Prikry construction. It therefore does not
answer whether a cover-exacting cardinal can lie above a strongly compact
cardinal. That is the explicit question in the 2026 notes.\cite[Problem~5.2]{cover}

It also does not conflict with the cover-exacting obstruction above an
extendible. If that obstruction is used, any model produced here has no
extendible below the cover-exacting witness. The positive construction is
about isolated existence, not the forbidden ordering of two cardinals.
A statement that the forcing preserves ``large cardinals'' without
specifying their position and preservation theorem would conceal exactly
the missing step.

Finally, there is no claim of a new equiconsistency for ultraexactingness.
The proof internalizes some branch of a tree; it does not show that the
chosen local embedding has its restriction to $V_\lambda$ in its domain.
The ultraexacting clause cannot be read into the conclusion.

\section{A definable-family inconsistency theorem}\label{sec:families}

\subsection{The only large-cardinal input}

For $\lambda$ exacting, the established regularity theorem says
\begin{equation}\label{eq:hodregularity}
 \cf^V(\lambda)=\omega,\qquad
 \HOD_{V_\lambda}\models\text{``$\lambda$ is a regular cardinal.''}
\end{equation}
The second assertion is Theorem~2.10 of
Aguilera--Bagaria--L\"ucke.\cite{abl}
We do not reproduce its exacting-embedding proof. Everything in the
present section after \eqref{eq:hodregularity} is an explicit deduction
from it.

Write $N=\HOD_{V_\lambda}$. A set in $\OD_{V_\lambda}$ need not belong
to $N$: its members might not be definable from permitted parameters.
For example, a definable set of cofinal maps may have no individually
definable member. On the other hand, an $\OD_{V_\lambda}$ \emph{set of
ordinals}, or an $\OD_{V_\lambda}$ graph between ordinals, does belong
to $N$. Its transitive closure adds only hereditarily definable finite
codes and ordinals. This observation is the bridge used below.

\begin{theorem}[Projection width]\label{thm:width}
Let $\lambda$ be exacting, $0<\rho<\lambda$, and let
$F\in\OD_{V_\lambda}$ be a nonempty set of cofinal maps
$f:\rho\to\lambda$. Define
\[
 P_\xi=\{f(\xi):f\in F\}\qquad(\xi<\rho).
\]
There is $\xi<\rho$ with $P_\xi$ unbounded in $\lambda$. For any such
$\xi$,
\begin{equation}\label{eq:fullwidth}
 \ot(P_\xi)=\lambda,\qquad |P_\xi|^V=\lambda,
 \qquad |F|^V\geq\lambda.
\end{equation}
If all members of $F$ are increasing, the first two equalities hold for
every coordinate greater than or equal to some $\xi<\rho$.
\end{theorem}

\begin{proof}
Fix one definition of $F$ using ordinal parameters and a finite tuple
$\vec x$ of members of $V_\lambda$. All definitions in this proof use
that same tuple, together with ordinals. This uniformity is important.
Each $P_\xi$ is a set of ordinals definable from these parameters, and
therefore belongs to $N$.

Suppose every $P_\xi$ were bounded in $\lambda$. Define
\[
 b(\xi)=\sup\{\beta+1:\beta\in P_\xi\}\quad(\xi<\rho).
\]
Each value is below $\lambda$. The \emph{entire graph} of $b$ is
definable from $\vec x$ and the fixed ordinal parameters by substituting
the definition of $F$. It is a graph between ordinals, so $b\in N$.
Since $\rho<\lambda$ and $\lambda$ is regular in $N$,
\[
 \sup_{\xi<\rho}b(\xi)<\lambda.
\]
Every $f\in F$ is pointwise bounded by $b$, contradicting the cofinality
of $f$. Thus some $P_\xi$ is unbounded.

The increasing enumeration of this $P_\xi$ is definable from $P_\xi$
and belongs to $N$. If its order type were below $\lambda$, it would
be a short cofinal map in $N$, again contradicting regularity. A subset
of the ordinal $\lambda$ has order type at most $\lambda$, so its order
type is exactly $\lambda$. This is an absolute order-type computation.
Since $\lambda$ is also an ambient cardinal, $|P_\xi|^V=\lambda$.
The evaluation map $F\to P_\xi$ is onto, so ambient Choice gives
$|F|^V\geq\lambda$.

Finally, suppose all $f\in F$ are increasing and $P_\xi$ is unbounded.
For every $\eta\geq\xi$ and every $\beta<\lambda$, choose $f\in F$
with $f(\xi)>\beta$. Then $f(\eta)\geq f(\xi)>\beta$. Hence
$P_\eta$ is unbounded as well, and the same argument applies.
\end{proof}

\begin{corollary}[Small-family incompatibility]\label{cor:smallfamily}
The following theory is inconsistent:
\[
\begin{split}
 \ZFC+\exists\lambda\exists\rho\exists F\,[\,&\Ex(\lambda)
 \wedge 0<\rho<\lambda\wedge F\in\OD_{V_\lambda}\wedge 0<|F|<\lambda\\
 &\wedge\ \forall f\in F\ (f:\rho\to\lambda\text{ is cofinal})\,].
\end{split}
\]
No definable selection of a member of $F$ is required.
\end{corollary}

\subsection{A local covering formulation}

\begin{definition}
Let $\mathsf{SmallODCover}(\lambda)$ assert that every
$x\in V_{\lambda+1}$ belongs to a set $A\in\OD_{V_\lambda}$ of ambient
cardinality less than $\lambda$.
\end{definition}

This is a definability covering principle, not the embedding-relative
cover-exacting property. In particular, it must not be substituted for
Definition~\ref{def:relative} merely because both use the word ``cover.''

\begin{corollary}\label{cor:odcover}
\[
 \ZFC\vdash
 \Ex(\lambda)\ \Longrightarrow\ \neg\mathsf{SmallODCover}(\lambda).
\]
More locally, if $c:\omega\to\lambda$ is cofinal and $\lambda$ is
exacting, every $A\in\OD_{V_\lambda}$ with $c\in A$ has
$|A|\geq\lambda$.
\end{corollary}

\begin{proof}
Given such an $A$, restrict it to
\[
 F=\{f\in A:f:\omega\to\lambda\text{ is cofinal}\}.
\]
This family is nonempty, belongs to $\OD_{V_\lambda}$, and is a subset
of $A$. Theorem~\ref{thm:width} implies $|A|\geq|F|\geq\lambda$.
Since \eqref{eq:hodregularity} supplies a cofinal $c$ of rank $\lambda$,
we have $c\in V_{\lambda+1}$ and the covering principle fails.
\end{proof}

\subsection{What this changes in the negative programme}

The supplied report sought a single cofinal map in HOD or
$\HOD_{V_\lambda}$.\cite[Section~5.4]{unified}
The new deduction shows that a smaller demand suffices: construct a
\emph{definable family of candidates} of size below $\lambda$.
No canonical least member and no definable well-order of the family are
necessary. Taking coordinate suprema extracts the contradiction without
making a choice inside the definability model.

This is nevertheless a conditional inconsistency theorem. We have not
shown that strong compactness or extendibility below an ordinary exacting
cardinal produces such a family. It would be incorrect to replace the
missing family-construction theorem by a choice of candidates under an
arbitrary ambient well-order.

The cloud construction does not violate this result. Its family may
have size exactly $\lambda$, and its definition uses the generic sequence
and ground-name parameters, not necessarily parameters in $V_\lambda$.
Both distinctions are substantive. The positive argument requires a
small invariant predicate; the negative argument requires a still smaller
family definable from the specified low-rank parameters.

\section{A complementary approximation obstruction}\label{sec:approximation}

The regularity gap also yields a useful diagnostic for proposed covering
or approximation arguments. This is an elementary fresh-sequence
consequence, not a separate large-cardinal lower-bound theorem.

\begin{proposition}[A fresh countable cofinal set]\label{prop:fresh}
Let $\lambda$ be exacting, let $N=\HOD_{V_\lambda}$, and choose an
increasing cofinal set $c\subseteq\lambda$ of order type $\omega$ in $V$.
Then $c\notin N$, but
\begin{equation}\label{eq:fresh}
 c\cap a\in N
\end{equation}
for every $a\in N$ for which $N$ contains an injection
$a\to\rho$ for some $\rho<\lambda$. In fact, each intersection in
\eqref{eq:fresh} is finite.
\end{proposition}

\begin{proof}
If $c\in N$, its increasing enumeration would be a countable cofinal
map in $N$, contrary to regularity of $\lambda$ there. Now take $a$
as in the statement. The set $a\cap\lambda$ belongs to $N$ and is
well-orderably smaller than $\lambda$ there, hence is bounded in
$\lambda$. An increasing cofinal sequence of order type $\omega$ has
only finitely many points below any fixed bound less than $\lambda$.
Thus $c\cap a$ is a finite set of ordinals and belongs to $N$.
\end{proof}

To avoid concealing a choice assumption, call the following the
\emph{well-orderable $\delta$-approximation test}: whenever $A\subseteq N$
is an ambient set and $A\cap a\in N$ for every $a\in N$ admitting
an injection in $N$ into some ordinal below $\delta$, then $A\in N$.
Proposition~\ref{prop:fresh} refutes this test for every infinite
$\delta\leq\lambda$, with the same witness $c$.
When the inner model satisfies Choice, this is the usual small-set
form of the approximation test. No blanket assertion of Choice for
$\HOD_{V_\lambda}$ is used here.

This explains a limitation on one possible attack from the original
report. An approximation principle strong enough to recover $c$ cannot
simply be assumed as a harmless background regularity property of the
HOD pair: in the exacting configuration it is already a contradictory
extra condition. A successful argument must derive that principle from
some other stated hypothesis, not from the existence of the outer
universe's cofinal sequence alone.

\section{Proof audit and the remaining research targets}\label{sec:audit}

\subsection{The interfaces checked in the construction}

\begingroup\small
\renewcommand{\arraystretch}{1.17}
\begin{longtable}{@{}P{.27\textwidth}P{.66\textwidth}@{}}
\toprule
\textbf{Potential failure}&\textbf{How it is addressed}\\\midrule\endfirsthead
\toprule\textbf{Potential failure}&\textbf{How it is addressed}\\\midrule\endhead
An ill-founded direct limit&Only the coded transitive set $J(N_0)$ is used. Its height is $J(\delta)$ for $\delta<\kappa^+$, exactly where the $\Iwf$ hypothesis guarantees well-foundedness.\\[4pt]
An arbitrary name is not fixed&Equation~\eqref{eq:fixedrange} fixes every $J(a)$ with $a\in N_0$. It does not require $a$ to be ordinal definable.\\[4pt]
The generic sequence is not fixed&It moves to its tail. The family is invariant under finite changes of its range, so the family, rather than its chosen member, is fixed.\\[4pt]
The predicate has the wrong rank&The name is forced to belong to $V_a$ over $N$. All finite-change generics generate the same $W$, so every evaluation lies in the same $V_a^W$.\\[4pt]
The cardinal bound varies with height&The family of finite modifications has size at most $\Lambda$ in $W$, independently of the name and of $a$. This is the transported version of the single bound $\kappa$.\\[4pt]
An external embedding is used as an internal witness&The functions $q_n$ have rank below $\Lambda$ and belong to $W$. They code high-rank partial maps; a tree in $W$ then has an internal branch by ordinal-rank absoluteness.\\[4pt]
An image set is confused with a pointwise image&$e(Y)=Y$ follows by elementarity of the defining formula. The proof does not assert $e``Y=Y$.\\[4pt]
A global scheme is assumed in a small rank&The collapse is taken from a rank reflecting a fixed finite collection of formulas. Appendix~\ref{app:fragment} identifies their roles.\\[4pt]
A set of definable candidates is treated as hereditarily definable&The width proof only puts the ordinal projections and their uniformly defined bound function into $\HOD_{V_\lambda}$. It never puts all of $F$ there.\\[4pt]
A consistency bound becomes equiconsistency&Only the displayed one-way reductions are claimed. No reverse arrow from cover exactingness to $\Iwf$ is established.\\
\bottomrule
\end{longtable}
\endgroup

\subsection{Three focused next questions}

The first question concerns the middle of the revised interval:
\begin{question}
Can the $\Iwf$ upper bound be lowered further for
$\exists\lambda\,\CEx_\lambda(\lambda)$, or can the cover property supply
an iteration lower bound not already implied by ordinary exactingness?
\end{question}
The construction identifies a specific possible obstruction to lowering
it: the chosen counterexample hull must remain well-founded after
iteration. One would need either a smaller well-foundedness certificate
for those particular hulls or a different internalization setup.

\begin{question}
Under a smaller strongly compact cardinal, can one construct an
$\OD_{V_\lambda}$ family of fewer than $\lambda$ candidates containing
one short cofinal map into $\lambda$?
\end{question}
An affirmative derivation from the exacting--strongly-compact hypotheses
would yield an inconsistency through Corollary~\ref{cor:smallfamily}.
The family bound must be \emph{strictly} below $\lambda$, and the allowed
parameters must be verified. This question is a sufficient local route,
not an assertion that the missing construction follows from compactness.

\begin{question}
Can a cover-exacting construction retain a specified strongly compact
cardinal below the cover-exacting witness, while avoiding the
club-definability obstruction proved for extendibles?
\end{question}
The present construction gives no such preservation theorem. Its
successful isolated-existence case is a control against arguments that
would refute cover exactingness without using the smaller cardinal.

The global cardinal-preserving and Icarus tracks are not resolved by
these results. In particular, nothing here turns internal stationarity
in a target inner model into ambient stationarity. The original report's
stationary-collision diagnostic still requires that separate input.

\subsection{Overall conclusion}

The continuation produces a positive consistency refinement and a
negative local criterion that work in opposite directions without
conflict. A Prikry name can be placed in a family of at most the critical
limit's cardinality whose defining class is invariant under the shift.
A low-rank-parameter-definable family of \emph{fewer} than that many
cofinal maps is impossible at an exacting cardinal.

\clearpage
\appendix
\section{Finite fragments, ranks, and the model ledger}\label{app:fragment}

\subsection{Why the small models have enough set theory}

A transitive collapse of a hull in $V_\zeta$ must not be called a model
of all of $\ZFC$ merely because $\zeta$ was chosen large. The proof uses
the following standard finite-fragment procedure instead.

First fix the formulas defining the forcing relation for
$\Bad(\kappa,\alpha,y)$ and its negation, the set-satisfaction relation,
Prikry-name evaluation, the finite-change family, and the tree of finite
partial maps. Include the finite collection of ZFC consequences used to
construct these objects and to carry out their proofs. In particular,
include the needed instances of well-founded recursion and of the
forcing theorem, normal-measure Prikry genericity, the elementary hull
construction, and the assertion that a well-founded set relation has
an ordinal rank.

Each of these is a uniform assertion or uses finitely many instances
in a fixed formal proof. There is therefore a finite subtheory of
$\ZFC$ sufficient for the argument. Enlarge the reflecting list to
contain this theory and the displayed forced-counterexample statement.
L\'evy reflection supplies a rank $V_\zeta$ satisfying and correctly
reflecting this finite list, with $\zeta$ above all specified parameters.
Its elementary hull and transitive collapse satisfy the same list.
The direct-limit map transfers it to the well-founded $N$, and its
included forcing theorems apply to $N[G_c]$.

This is a finite-proof reflection argument, not a demand for
$V_\zeta\prec V$ for all formulas simultaneously. The internalization
lemma quantifies over all first-order formulas of \emph{set structures}
using their single set-satisfaction relation; it does not require a
truth predicate for the universe.

\subsection{The universe changes that must not be conflated}

\begin{center}
\small
\renewcommand{\arraystretch}{1.18}
\begin{tabularx}{\textwidth}{@{}lY@{}}
\toprule
\textbf{Model}&\textbf{Role}\\\midrule
$V$&Original ambient universe, containing $j$ and the whole construction of the local auxiliary model.\\
$V[G]$&Hypothetical Prikry extension in which failure at $\kappa$ is assumed, in order to obtain the forced counterexample.\\
$N_0$&Small transitive collapse of a hull of a reflecting $V_\zeta$; its cardinality is $\kappa$ and its ordinal height is below $\kappa^+$.\\
$M^j_\omega$&Possibly ill-founded direct limit. It is used as a transport structure, not as the final universe.\\
$N=J(N_0)$&Well-founded, transitive, sufficiently correct small model containing the ambient $V_\Lambda$.\\
$W=N[G_c]$&Auxiliary model inside the original $V$. Its external shift is internalized only through a tree argument. It is not $V[G]$.\\
$M[H]$&Countable transitive model of the final theory obtained at the end of Corollary~\ref{cor:transitive}.\\
\bottomrule
\end{tabularx}
\end{center}

The counterexample is transported by
\[
 V\models\mathbf1\forces\Bad(\kappa,\alpha,\tau)
 \ \Longrightarrow\
 N\models\mathbf1\forces\Bad(\Lambda,a,t)
 \ \Longrightarrow\
 W\models\Bad(\Lambda,a,t^{G_c}).
\]
The contradiction is entirely \emph{inside $W$}. It then rules out the
initial forced failure and proves the result about $V[G]$.

\subsection{Rank checks}

The membership relation $E_0$ on $\kappa$ is a subset of
$\kappa\times\kappa$ and belongs to $V_{\kappa+1}$, using the usual
finite-rank coding of ordered pairs. The coding bijection and the
transitive collapse have rank below $\kappa^+<\Lambda$.
The image relation $E$ is a subset of $\Lambda\times\Lambda$, with
rank at most $\Lambda$, so it belongs to $V_{\Lambda+1}$.

For internalization, $q_n=e\restr\kappa_n$ is a subset of
$\kappa_n\times\kappa_{n+1}$. Its rank is below $\Lambda$, so
$q_n\in W$ follows from rank inclusion, not from a closure principle
for $W$ at length $\Lambda$. The actual partial map $p_n$ may have
much higher rank, since its domain lies in $V_a^W$. It is nevertheless
in $W$ because it is constructed there from $q_n,f,g$.

A cofinal map $c:\omega\to\lambda$ has rank $\lambda$ and belongs
to $V_{\lambda+1}$, not $V_\lambda$. This is why its use as a parameter
in the cloud definition is not permitted in the negative theorem's
$\OD_{V_\lambda}$ family hypothesis.

\section{Source and verification record}\label{app:verification}

The supplied \texttt{Large\_Cardinals\_Unified\_Report.tex} was read as
the starting research document. Its palette, font packages, page geometry,
heading styles, theorem styling, and assessment-box design are retained
here. The new text is a continuation rather than an alteration of the
original report's claims.

The primary sources checked for the main argument are the July 2026
cover-exacting notes, version 1 of \emph{Large cardinals beyond HOD},
version 2 of the iterability paper, and version 4 of the original
exacting-cardinal paper. The cover notes explicitly acknowledge that
related consistency constructions yield cover variants; this is why
historical novelty is not inferred from the absence of the exact
$\Iwf$ statement in the supplied report.

The proof audit concentrated on the fixedness of arbitrary transported
name parameters, the finite-change equivalence relation, the internal
cardinality bound, the well-founded part of the iterate, and the internal
existence of the final embedding. No finite computation is offered as
evidence for the infinite set-theoretic assertions. The manuscript has
not been formalized in Lean or another proof assistant, and it has not
undergone independent specialist review.

The archive includes the compilable source and rendered PDF, with a
build script and a short dependency ledger. No external figure assets or
font files are needed in the archive. The original Palatino-style text
and mathematical typography are selected by the same \texttt{newpxtext}
and \texttt{newpxmath} packages as the supplied source.

\clearpage
\phantomsection
\addcontentsline{toc}{section}{References}
\begin{thebibliography}{99}
\small
\setlength{\itemsep}{7pt}
\setlength{\parskip}{1pt}

\bibitem{unified}
\emph{Large cardinals at the inconsistency frontier: a unified assessment
of ZFC-compatible hypotheses, their obstructions, and the evidence on
both sides}. Supplied research report prepared for Vladimir,
18 September 2026. Source file:
\nolinkurl{Large_Cardinals_Unified_Report.tex}.
The starting document, particularly its sections on exactingness,
cover exactingness, and the proposed research programme.

\bibitem{cover}
G.~Goldberg. \emph{Consistency beyond the Kunen inconsistency}.
Lecture notes dated 1 July 2026, reporting joint work with D.~Blue.
\url{https://math.berkeley.edu/~goldberg/Slides/CoverExact.pdf}.
Definitions~2.6 and~3.1; Theorem~3.5 and the subsequent remark;
Theorem~3.6; Problem~5.2. Used as lecture notes, not described as a
refereed publication.

\bibitem{abgl}
J.~P.~Aguilera, J.~Bagaria, G.~Goldberg, and P.~L\"ucke.
\emph{Large cardinals beyond HOD}.
\arxiv{2509.10254v1}, 12 September 2025.
Proposition~5.1, Definition~5.2, the proof of Theorem~5.4,
Corollary~5.5, Theorem~5.10, and Corollary~5.11 are the principal
locators for the iteration-localization and consistency comparisons.

\bibitem{ad}
A.~Andretta and V.~Dimonte. \emph{The iterability hierarchy above I3}.
Archive for Mathematical Logic \textbf{58} (2019), 77--97.
\doi{10.1007/s00153-018-0624-5}; \arxiv{1712.03877v2}.
Section~3, especially Lemmas~3.3--3.5, supplies the iteration and
amenable-predicate elementarity background.

\bibitem{abl}
J.~P.~Aguilera, J.~Bagaria, and P.~L\"ucke.
\emph{Large cardinals, structural reflection, and the HOD Conjecture}.
\arxiv{2411.11568v4}, 16 September 2025.
Theorem~2.10 supplies regularity of an exacting cardinal in
$\HOD_{V_\lambda}$; the conventions immediately preceding it specify
$\OD_\Gamma$ and $\HOD_\Gamma$.

\end{thebibliography}
\end{document}
